% Transcription — fonds Alexandre Grothendieck, shelfmark 135, batch 1 (pages 1-20).
% Established from the facsimile archives/batches/135/batch-01.pdf (a mirror of
% https://grothendieck.umontpellier.fr/135.pdf), per the skill
% transcribe-grothendieck. Pages 1 and 6 are blank — absent. Pages 3 and 5 are
% two copies of the same English typescript, a song titled "Lonely Road" (the
% second copy struck through twice in ink) — reused paper, not the mathematics
% in hand, absent. The folder belongs to the inventory group 134-1-8 (Pursuing
% stacks); the notes are on Gr-categories, Picard categories, and their
% classification by 2-gerbes — the circle of ideas of Hoàng Xuân Sính's thesis
% and of Breen's later work, consistent with the inventory's dating.
% Pass: Fable 5 (claude-fable-5), 2026-08-09 — first pass, unchecked against the pages by a human.
\documentclass[11pt,a4paper]{article}
% Le préambule commun à toutes les transcriptions du fonds Grothendieck.
%
% Il tient en une idée : **l'incertitude se compose, elle ne se résout pas.**
% Une transcription qui lisse un mot douteux a détruit la seule chose pour
% laquelle on la faisait — un lecteur doit pouvoir savoir, à chaque endroit,
% ce qui était sur le feuillet et ce qui a été ajouté. Les six macros qui
% suivent sont donc l'appareil critique, et non de la décoration.
%
% Les mêmes commandes sont reconnues par `scripts/render.mjs`, qui produit la
% vue de lecture du site. Une macro ajoutée ici doit l'être aussi là-bas,
% faute de quoi le rendu échouera bruyamment — ce qui est voulu : un rendu
% silencieusement faux est pire qu'un rendu absent.

\ProvidesPackage{grothendieck}[2026/08/08 Transcriptions du fonds Grothendieck]

\usepackage{amsmath,amssymb,amsthm}
\usepackage{mathrsfs}
\usepackage{stmaryrd}
% Les diagrammes commutatifs sont partout dans ce fonds : c'est la langue
% même de la géométrie algébrique de Grothendieck, et une transcription qui
% les rendrait en prose ne transcrirait rien.
\usepackage{tikz-cd}
% Français par défaut — c'est la langue des feuillets — mais l'anglais est
% chargé aussi : la traduction et le résumé sont des documents anglais, et la
% typographie française (espaces avant les deux-points) y serait une faute.
% Ces documents basculent par \selectlanguage{english} en tête de corps.
\usepackage[english,french]{babel}
\usepackage{fontspec}
\usepackage{geometry}
\geometry{a4paper,margin=25mm,marginparwidth=28mm}
\usepackage{marginnote}
\usepackage{xcolor}
% ulem, not soul. `soul`'s \st and \ul are built on a character-by-character
% scan that XeLaTeX's font machinery breaks: the document fails with an
% undefined control sequence rather than degrading. ulem does the same two jobs
% and is Unicode-engine safe. `normalem` stops it hijacking \emph, which the
% transcriptions use for Grothendieck's own emphasis.
\usepackage[normalem]{ulem}
\usepackage{hyperref}
\hypersetup{colorlinks=true,linkcolor=black,urlcolor=black,pdfborder={0 0 0}}

% --- Métadonnées du lot -----------------------------------------------------
% Reprises telles quelles par le script de rendu, qui les affiche en tête de
% la vue de lecture. Elles fixent ce que le fichier prétend couvrir : sans
% elles, une transcription flotte sans point d'attache dans un fonds de seize
% mille feuillets.

\newcommand{\folder}[1]{\gdef\grothFolder{#1}}
\newcommand{\batch}[1]{\gdef\grothBatch{#1}}
\newcommand{\leaves}[2]{\gdef\grothFirst{#1}\gdef\grothLast{#2}}
\newcommand{\foldertitle}[1]{\gdef\grothFoldertitle{#1}}
\gdef\grothFolder{?}\gdef\grothBatch{?}\gdef\grothFirst{?}\gdef\grothLast{?}\gdef\grothFoldertitle{}

% --- Le résumé --------------------------------------------------------------
%
% Il ouvre la lecture modernisée, sous le titre « Résumé », et il est composé
% en petit. Ce n'est pas un ornement : c'est le seul passage du document qui
% ne soit pas des mathématiques, et un lecteur doit pouvoir le reconnaître
% avant de l'avoir lu — puis le sauter s'il connaît déjà le sujet, ou s'y
% arrêter s'il ne le connaît pas. Le filet qui le suit marque la reprise du
% corps.

\newenvironment{resume}
  {\small\setlength{\parskip}{0.45em}}
  {\par\medskip\hrule\medskip}

% --- Mots-clés ---------------------------------------------------------------
%
% En anglais, en clôture du résumé de la lecture modernisée : le vocabulaire
% moderne sous lequel chercher ce que les feuillets construisent. Le manifeste
% du site (`npm run manifest`) les extrait pour étiqueter la cote entière —
% cette ligne est la source unique des tags, il n'y a pas de fichier à côté.
\newcommand{\keywords}[1]{\par\smallskip\noindent
  {\footnotesize\textsc{Keywords} --- \textit{#1}}\par}

% --- L'appareil critique ----------------------------------------------------

% Le numéro de feuillet, en marge. C'est l'ancre qui permet de régler un
% désaccord sur une formule en regardant *un* feuillet plutôt qu'une chemise
% de six cents. Le site s'en sert aussi pour tourner les pages du fac-similé
% quand on fait défiler la transcription.
\newcommand{\leaf}[1]{%
  \marginnote{\footnotesize\color{gray!70}\textsf{#1}}%
  \label{leaf:#1}\ignorespaces}

% Illisible : on ne devine pas. Un mot inventé qui se lit comme les autres est
% le pire résultat possible.
\newcommand{\ill}{\textcolor{red!60!black}{[\,\ldots\,]}}

% Lecture douteuse : le mot est donné, et signalé comme incertain.
\newcommand{\uncertain}[1]{\uline{#1}}

% Ajout de l'éditeur — une lettre manquante, un mot sous-entendu.
\newcommand{\add}[1]{\textcolor{blue!50!black}{[#1]}}

% Biffé par Grothendieck lui-même. Ce qu'il a rayé fait partie du document :
% une rature dit souvent où la pensée a changé de direction.
\newcommand{\struck}[1]{\sout{#1}}

% Note du transcripteur, sur l'état matériel du feuillet.
\newcommand{\note}[1]{\par\smallskip{\small\color{gray!60!black}\textit{[#1]}}\par\smallskip}

% Note marginale de Grothendieck, distincte du corps du texte.
\newcommand{\marginal}[1]{\par\smallskip\noindent\hspace{1em}%
  {\small\color{gray!60!black}#1}\par\smallskip}

% --- Titre ------------------------------------------------------------------

\AtBeginDocument{%
  \begin{center}
    {\large\bfseries Fonds Alexandre Grothendieck --- cote n\textsuperscript{o}~\grothFolder}\\[2pt]
    {\itshape\grothFoldertitle}\\[4pt]
    {\small feuillets \grothFirst--\grothLast\ (lot \grothBatch)}\\[2pt]
    {\footnotesize\color{gray!60!black}Transcription établie d'après le fac-similé numérisé
     par l'Université de Montpellier.\\
     Les lectures incertaines sont soulignées, les illisibles notées [\,\ldots\,],
     les ajouts entre crochets.}
  \end{center}
  \vspace{1em}}

\endinput


\folder{135}
\batch{1}
\pages{1}{20}
\dating{[vers 1974]}
\watermark{Édition de démonstration}
\foldertitle{Catégories [et Gr-catégories] : notes manuscrites (s.d.), copies de tapuscrit annoté (s.d.), copies de manuscrit annoté (s.d.).}

\begin{document}

\section*{« Tordre » une catégorie de Picard (pages 2 et 4)}

\page{2}
\uncertain{« Tordre »} une catégorie de Picard

1) Multiplier une Gr-catégorie par un groupe.

Soit $C$ une (gr)-cat, $\Gamma$ un groupe — on définit une nouvelle
(Gr)-cat.\ notée $\Gamma \times C$. Le plus simple est de dire que $\Gamma$
désigne une catégorie discrète \uncertain{d'ensemble sous-jacent} $\Gamma$,
et la loi de composition de $\Gamma$ définit un foncteur $\otimes$, avec une
\struck{structure} \uncertain{contrainte} d'associativité stricte
\struck{évidente}. [Il \uncertain{existe} une contrainte de commutativité
ssi $\Gamma$ est commutatif. (Les deux \uncertain{contraintes} sont alors
compatibles, et $\Gamma$ est une catégorie \uncertain{ACU}.)] On
\uncertain{retrouve ainsi} les produits des (Gr-cat) $\Gamma$
\uncertain{et} $C$.
\[
\pi_0(\Gamma \times C) = \pi_0(\Gamma) \times \pi_0(C)
= \Gamma \times \pi_0(C) \xrightarrow{\ p\ } \pi_0(C)
\]
\[
\pi_1(\Gamma \times C) = \pi_1(\Gamma) \times \pi_1(C) = \pi_1(C),
\qquad
k_{\Gamma \times C} = p^{*}(k_C).
\]

2) Supposons que $C$ soit = une Pic-cat, et \uncertain{$\Gamma$}
commutatif. \struck{Alors} $\Gamma \times C$ est une Pic-catégorie.
\note{le signe « = » tient lieu de « même » sous la plume de l'auteur ;
l'abréviation revient tout au long du dossier}

3) Si $C$ est une Pic-catégorie, les \uncertain{autres} contraintes de
comm.\ qui en font (avec sa contrainte d'associativité) une Pic-catégorie
\ill{} compatibles avec la contrainte d'ass., correspondent aux
applications bilinéaires antisymétriques
\[
\varphi : \pi_0(C) \times \pi_0(C) \longrightarrow \pi_1(C).
\]
L'application correspondante $\pi_0(C) \xrightarrow{\ \Sigma_\varphi\ }
\pi_1(C)$ est donnée par $\Sigma_\varphi(x) = \varphi(x,x)$.

\page{4}
4) \textbf{Lemme.} Soient $M$, $N$ deux groupes commutatifs, et
$\Sigma : M \to {}_2N$ un hom. Alors il existe une application bilinéaire
antisymétrique $\varphi : M \times M \to N$ telle que
$\varphi(x,x) = \Sigma(x)$ pour tout $x \in M$.

Remplaçant $M$, $N$ resp.\ par \uncertain{$M_2$}, ${}_2N$, on
\uncertain{peut} supposer $M$, $N$ des vectoriels sur $\mathbf{F}_2$.
\note{la lettre qui précède $M_2$ est indistincte ; comprendre sans doute
$M/2M$}
Soit $(e_i)_{i \in I}$ une base de $M$ ; \uncertain{alors $\Sigma$ est
donnée} par les $n_i = \Sigma(e_i) \in N$, \uncertain{tandis que}
$\varphi$ \uncertain{sera} donnée par les $\varphi(e_i, e_j) = n_{ij}$, la
condition d'antisymétrie (= \uncertain{symétrie}) étant $n_{ij} = n_{ji}$.
On prendra
\[
\varphi(e_i, e_j) =
\begin{cases}
0 & \text{si } i \neq j \\
n_i & \text{si } i = j
\end{cases}
\]
\uncertain{OK}.

\textbf{Corollaire.} Toute catégorie de Picard devient stricte en
changeant sa seule contrainte de commutativité. \uncertain{En tant que}
Gr-catégorie, elle est \uncertain{de type triviale}\ldots
\note{les deux derniers membres de phrase sont des ajouts interlinéaires
serrés}

5) Soit $C$ une catégorie de Picard, $\Gamma$ un groupe commutatif,
$\varphi : \Gamma \times \Gamma \to \pi_1(C)$ une application bilinéaire
antisymétrique. On désigne par $\Gamma \times^{\varphi} C$ la catégorie de
Picard déduite de $\Gamma \times C$ (cf n°2) par l'application bilinéaire
antisym.
\[
\overline{\varphi} : \pi_0(\Gamma \times C) \times \pi_0(\Gamma \times C)
\longrightarrow \pi_1(\Gamma \times C)
\]
(cf n°3) donnée par $\overline{\varphi}((\alpha,x),(\beta,y)) =
\varphi(\alpha,\beta)$. \uncertain{Avec comparaison de}
\[
\pi_0(\Gamma \times^{\varphi} C) \approx \pi_0(\Gamma \times C),
\qquad
\pi_1(\Gamma \times^{\varphi} C) \cong \pi_1(C).
\]
\note{le premier membre de gauche ne laisse pas voir le $\varphi$ ; sous
ces formules, un griffonnage à demi biffé — une flèche vers
« $F \times \Gamma$ » et deux flèches marquées $\nu$?, $\varphi$? — que
nous ne savons pas lire}

6)
\note{le numéro reste seul : la note s'interrompt là}

\section*{Représentations d'un complexe semi-simplicial et 2-gerbes sur
$B_G$ (pages 7 à 10)}

\page{7}
\note{page 1 de la main de l'auteur}
1. $K$ complexe semi-simplicial, $P$ \struck{Gr-catégorie}
$\otimes$-catégorie \uncertain{AU}. « Représentation \uncertain{$\Lambda$}
de $K$ dans $P$ » : données

a) $\forall\, u \in K_1$, $\Lambda(u) \in P$ ;

b) $\forall\, \sigma \in K_2$, posant $u = d_1\sigma$, $v = d_3(\sigma)$,
$w = d_2(\sigma)$ :
\[
\Lambda(u) \otimes \Lambda(w) \simeq \Lambda(v) ;
\]
\note{les lettres $v$ et $w$ se confondent sur la page, et le premier et le
dernier indice de face sont douteux ; avec les indices tels que lus, la
composition demanderait plutôt $\Lambda(u) \otimes \Lambda(v) \simeq
\Lambda(w)$}

c) $\forall\, \Delta \in K_3$, \ldots{} « axiome des
\uncertain{pentagones} ».

La composée $K_1 \to P \to \pi_0(P)$ définit un « homomorphisme »
$K_1 \xrightarrow{\ \alpha\ } \pi_0(P) = G$, d'où un « homomorphisme »
$K_1 \xrightarrow{\ \overline{\alpha}\ } \mathrm{Aut}(\pi_1(P))$
[$\pi_1(P) = M$], d'où un système local sur $K$ de valeur $M$, noté
$(M, \overline{\alpha}) = \overline{M}$.

\textbf{Pb} : quand $K$, $P$, $\alpha$ (donc $\overline{M}$) sont fixés, on
veut « classer » les représentations de $K$ dans $P$ correspondantes.

\textbf{Résultat} : ces représentations forment une 2-catégorie, qui est un
« 2-torseur » sous une 2-catégorie de Picard stricte d'invariants
$H^2(K, \overline{M})$, $H^1(K, \overline{M})$, $H^0(K, \overline{M})$.
\note{les arguments des trois $H$ portent des résidus biffés — l'auteur a
d'abord écrit $(M, \alpha)$ avant de contracter en $\overline{M}$}
Cette 2-catégorie est contravariante p.r.\ à $K$ variable, covariante
p.r.\ à $P$ variable.

Cas particuliers utiles pour la suite : $K$ est le nerf d'une catégorie
$\mathcal{X}$.

\marginal{$\mathrm{REP}(K ; P, \alpha)$, $\mathrm{REP}(\mathcal{X} ; P,
\alpha)$}

2. $P$ Gr-catégorie donnée, d'invariants $\pi_0(P) = G$, $\pi_1(P) = M$. On
va lui associer une 2-gerbe $\widetilde{P}$ sur $B_G$ — donc à
\uncertain{tout} $X \in \mathrm{Ob}\,B_G$ (ensemble où $G$ opère) on
\uncertain{va} associer une catégorie $\widetilde{P}(X)$, ainsi : soit
$\mathcal{X}$ le nerf de la catégorie \uncertain{sous-jacente à} $X$

\page{8}
($\mathrm{Ob}\,\mathcal{X} = X$, $\mathrm{Fl}\,\mathcal{X} = G \times X$
\ldots), avec l'application composée $K_1 = \mathrm{Fl}\,\mathcal{X} =
G \times X \xrightarrow{\ \mathrm{pr}\ } G$. On pose
\[
\widetilde{P}(X) = \mathrm{Rep}(\mathcal{X} ; P, \alpha_X)
= \mathrm{Rep}(X, P, \alpha_X)
\]
avec la contravariance \uncertain{p.r.\ à} $X$ évidente. On trouve que
c'est un 2-champ, qui est une 2-gerbe liée par $M_{(B_G)}$ (définie par
l'opération de $G$ sur $M$). \note{un mot biffé — « construit » ? — précède
la parenthèse}

Variante — $\widetilde{P}$ est covariante en $P$ dans un sens évident\ldots{}
En fait \uncertain{cela définit} un morphisme de la 3-catégorie des \ill{}
dans la 3-catégorie des 2-gerbes sur des $B_G$ variables.
\note{passage chargé de ratures ; un « Fixons » est biffé en cours de
phrase, et la marge gauche porte deux numéros d'alinéa dont l'un est biffé}

3. On a ainsi une équivalence de 3-catégories. Plus explicitement, $G$ et
le $G$-module $M$ étant fixés — \ldots{} une équivalence entre la
3-catégorie des Gr-catégories \struck{de Pic\ldots} épinglées par $(G,M)$,
et la 3-catégorie des 2-gerbes sur $B_G$ liées par $M_{(B_G)}$.

Soient $P$ Gr-catégorie épinglée par $G$, $M$ ; $P'$ —— $G'$, $M'$ ;
$G \xrightarrow{\ u\ } G'$, $M \xrightarrow{\ v\ } M'$ homs compatibles
avec \uncertain{les opérations}. Alors les $\otimes$-hom $P \to P'$
compatibles avec $G \to G'$, $M \to M'$ forment une 2-catégorie\ldots
\note{« épinglée par » corrige un « d'invariants » biffé ; la ligne de $P'$
répète celle de $P$ par un long tiret}

\page{9}
\note{page 2 de la main de l'auteur}
\ldots{}qui est un 2-\uncertain{ps-}torseur sous une 2-Pic-catégorie, et ce
2-torseur est 2-équivalent à \uncertain{celui des hom de 2-gerbes sur $B_G$
liées par $M'_{B_G}$, de $\widetilde{P} \otimes_{M_{B_G}} M'_{B_G}$ dans
$B_u^{*}(\widetilde{P}')$}. Les invariants de la 2-Pic-catégorie pertinente
sont $H^2(G,M')$, $H^1(G,M')$, $H^0(G,M')$. L'obstruction à trouver
\uncertain{un objet dans le} 2-ps-torseur est dans $H^3(G,M')$, et cela se
calcule comme différence de deux classes, images canoniques
\uncertain{d'éléments qui dérivent des} $H^3(G,M)$ et $H^3(G',M')$.

\marginal{$M_{B_G} \to M'_{B_G}$ ; $B_G \xrightarrow{\ B_u\ }
B_{G'}$}

4. Généralisation aux Gr-champs sur un topos $S$. \struck{Bien sûr,} si $G$
(groupe sur $S$) et $M$ ($G$-Module sur $S$) sont donnés, on interprète
tout en termes de 2-gerbes « relatives » sur $B_G$ \uncertain{mod} $S$,
liées par $M_{B_G}$, i.e.\ 2-gerbes sur $B_G$ trivialisées sur $S$ (par
$S \simeq B_e \to B_G$). Il faudrait \uncertain{la} mettre « en scène »
avec l'histoire des 1-gerbes relatives sur $B_G$, correspondant aux
extensions de $G$ par $M$. Il faudrait \uncertain{traiter} en = temps le
cas d'un $M$ non commutatif\ldots

\page{10}
5. Relations avec les « noyaux » de Mac Lane, ou « liens ». Soit $G$
groupe, $F$ groupe,
\[
\varphi : G \longrightarrow \mathrm{Autext}(F)
\]
hom (\uncertain{cela} $=$ se donner un lien sur $B_G$). On va lui associer
une Gr-catégorie \struck{de Picard} $P(\varphi)$ d'invariants $G$ et
$M = Z(F)$ (centre).

a) Description directe : on a la Gr-catégorie $P(F)$ des bitorseurs sous
$F$, qui a comme invariants
\[
\pi_0(P(F)) \in \mathrm{Autext}\,F ; \qquad \pi_1(P(F)) = Z(F),
\ \text{centre de } F.
\]
\note{le premier signe est bien « $\in$ » sur la page ; on attendrait
« $=$ »}
On définit \uncertain{donc} $P(\varphi)$ comme image inverse.

\struck{b) Description en termes des 2-gerbes liées par $M$.}

\textbf{NB.} Les trivialisations de $P(\varphi)$ correspondent aux
réalisations de $\varphi$ par une extension de $G$ par $F$
\uncertain{subordonnée} à $\varphi$. Il faudrait encore préciser ce point
en termes d'une équivalence de 2-catégories (2-torseurs sous une
2-catégorie de Picard). On retrouve ainsi le \uncertain{cas} de Mac Lane
dans $H^3(G,M)$, obstruction à la réalisabilité du « noyau » $\varphi$ par
une extension (cas \uncertain{analogue} du lien $\varphi$ par une
1-gerbe). \struck{Mais} le fait que \uncertain{tous les éléments} des
$H^3(G,M)$ (i.e.\ \uncertain{les catégories de Picard épinglées par
$(G,M)$}) puissent ainsi s'obtenir à l'aide de « noyaux » ne fait pas
\uncertain{(pour l'instant) l'objet d'une logique générale}.

\section*{Catégories et champs de Picard épinglés (pages 11 à 17)}

\page{11}
\note{page 3 de la main de l'auteur}
c) Soient $P$, $P'$ deux catégories de Picard, épinglées par $M$, $N$. Pour
qu'\uncertain{elles} soient équivalentes (avec leurs épinglages) il f.\ et
s.\ que les applications \uncertain{$\sigma_P$, $\sigma_{P'}$} $: M \to
{}_2N$ soient les =s. \uncertain{La catégorie des équivalences} est alors
\uncertain{un 1-torseur sous la} catégorie de \uncertain{Picard stricte}
\marginal{d'invariants $\mathrm{Ext}^1(M,N)$ et $\mathrm{Hom}(M,N)$
(obtenus \uncertain{par} $\mathrm{RHom}(M,N)$ \uncertain{tronqué})}

d) La classification\ldots{} des cat.\ de Picard épinglées par $M$, $N$,
\uncertain{à équivalence près}, est (\uncertain{comme groupe}) $\simeq
\mathrm{Hom}(M, {}_2N)$.

e) Revenant à la situation de c) et utilisant la multiplication de Baer des
Pic-cat.\ épinglées par $(M,N)$ — trouver que celles qui correspondent à un
$\sigma : M \to {}_2N$ donné forment une 2-catégorie qui est un 2-torseur
sous la 2-catégorie de Picard stricte formée des cat.\ de Pic.\ strictes
épinglées par $M$, $N$ [qui est un 2-groupoïde connexe : tous les objets
sont 2-équivalents\ldots]

f) Les Pic-catégories forment une \uncertain{2-catégorie AUC} de façon
évidente \uncertain{(et sur les objets et flèches)} \ldots
\note{la fin de la page, très raturée, s'interrompt sur « $D$, » ; la page
suivante reprend le sujet sous un nouveau numéro}

\page{12}
6. Pic-catégories. On peut définir les Pic-catégories comme \uncertain{les
objets} d'une 3-catégorie, qui \uncertain{sera} une sous-3-catégorie
(pleine au niveau des 2-objets et 3-objets) de celle des Gr-catégories.
\note{suit un passage biffé : « Soient $P$, $P'$ deux Pic-catégories, et
$\Lambda, \Lambda' : P \rightrightarrows P'$ deux $\otimes$-foncteurs
\ldots{} »}

Mais on peut aussi \uncertain{penser à} la 2-catégorie associée en gardant
les =s 0 et 1 objets, mais en prenant les 2-objets à isom.\ près (c'est en
fait la solution \ill{}) et \uncertain{oubliant} les 3-objets.
\uncertain{Si on se restreint} aux cat.\ de Picard strictes,
\uncertain{retrouve-t-on} alors la 2-catégorie de Deligne (où les
catégories Hom sont elles-\uncertain{mêmes} des catégories de Picard,
déduites des $\mathrm{RHom}(L_{.}, L'_{.})$ par troncature).

Si on s'intéresse \uncertain{surtout} aux Pic-cat.\ strictes, on peut dire
ceci :

a) Soit $P$ une Gr-catégorie. Elle admet une structure de Pic-catégorie
(\uncertain{une contrainte de comm.} compatible avec \ill{}) ssi elle est
splittable, i.e.\ $k(P) \in H^3(\pi_0, \pi_1)$ est \uncertain{nul}. Alors
$P$ admet = une structure de Picard stricte.

b) Les structures de Picard \struck{strictes} sur $P$ forment un torseur
sous \uncertain{$\mathrm{Bilant}$}$(\pi_0, \pi_0, \pi_1)$. Les structures
de Picard strictes forment un sous-torseur sous $\mathrm{Bilalt}(\pi_0,
\pi_0 ; \pi_1)$.
\note{le premier mot, surchargé d'une correction, se lit « Bilant » ou
« Bilalt » ; le second est nettement « Bilalt »}

\page{13}
\note{page 4 de la main de l'auteur}
\struck{\ill} \uncertain{Pour l'}addition ordinaire \uncertain{ils forment}
des triples $(M, N, \sigma : M \to {}_2N)$. Si $P$, $P'$ sont associés à
$(M,N,\sigma)$, $(M',N',\sigma')$, alors deux foncteurs $\Lambda, \Lambda'
: P \to P'$ \uncertain{$\otimes$-AUC} — ils définissent le = hom
$(M,N,\sigma) \to (M',N',\sigma')$ \uncertain{(dans la catégorie déduite de
(Pic) en passant au quotient par les 2-flèches)}. Mais bien sûr, la
catégorie des $\otimes$-foncteurs AUC $P \to P'$ qui sont subordonnés à un
hom donné $(M,N,\sigma) \to (M',N',\sigma')$ forment un groupoïde qui est
un torseur sous la Pic-catégorie définie par \uncertain{troncation} à
partir de $\mathrm{RHom}(M,N')$ (d'invariants $\mathrm{Ext}^1(M,N')$,
$\mathrm{Hom}(M,N')$).

\marginal{NB : $\mathrm{Hom}_{\otimes \mathrm{AUC}}(P, P')$ est un
\uncertain{groupoïde}}

7) \struck{Traductions} Extension aux champs de Picard sur un topos. Ils
forment une 2-$\otimes$-catégorie AUC, et définissent un 2-$\otimes$-champ
AUC. Il s'envoie dans la 1-$\otimes$-cat.\ AUC (ou 1-$\otimes$-champ AUC)
des $(M, N, \sigma : M \to {}_2N)$, mais si $\Lambda, \Lambda' : P \to P'$
définissent le = $(M,N,\sigma) \to (M',N',\sigma')$, p.ex.\ l'identité,
on\ldots

\page{14}
\note{page très raturée ; plusieurs raccords restent conjecturaux}
\ldots{}ne peut plus affirmer que $\Lambda$ et $\Lambda'$ sont isomorphes
(ni = loc.\ isom.), c'est tjrs vrai \struck{\ill} \uncertain{$\Lambda : P
\to P'$ fixés} si $\mathrm{Ext}^1(S; M, N') = 0$ (resp.\
$\mathrm{Ext}^1(M,N') = 0$). On ne peut affirmer non plus qu'un hom
\uncertain{$\Lambda$} de $(M,N,\sigma)$ en $(M',N',\sigma')$ \ldots{}
[c'est OK si $\mathrm{Ext}^2(S; M, N') = 0$, resp.\
\uncertain{$\mathrm{Ext}^2$}$(M,N') = 0$] \ldots{} ni que \uncertain{tout}
$(M,N,\sigma)$ provienne (\uncertain{i.e.\ d'un} $P$ — cf.\ plus bas —) ;
mais pour $P$, $P'$ \uncertain{$\Lambda$} fixés, le champ des \struck{cat.\
des} \uncertain{foncteurs AUC} $P \to P'$ correspondants est un 1-torseur
sous le champ de Picard strict d'invariants $\mathrm{Ext}^1(M,N')$,
$\mathrm{Hom}(M,N')$ obtenu par \uncertain{troncage} de
$\mathrm{RHom}(M,N')$ ; l'obstruction à trouver un tel
\uncertain{$\widehat{\Lambda}$} est dans $\mathrm{Ext}^2(S, M, N')$ [et il
\uncertain{devient nul} loc., et \uncertain{donc est} dans $H^0(S,
\underline{\mathrm{Ext}}^2(M,N'))$]. Tout ceci provient \uncertain{des jeux
habituels} : le Baer, et \uncertain{du fait qui en résulte}, que pour $M$,
$N$, $\sigma$ fixés, les champs de Picard d'invariants\ldots

\marginal{On peut \uncertain{prendre pour $\Lambda$} l'identité \ill}

\page{15}
\note{page 5 de la main de l'auteur}
$M$, $N$, $\sigma$ forment une 2-catégorie \struck{2-équivalente}
\uncertain{qui est un 2-ps-torseur} sous la catégorie \uncertain{des
champs} de Picard stricts d'inv.\ $M$, $N$ (et $\sigma$
\uncertain{donné}\ldots) — la « différence » de \struck{deux} $P$ et $P'$
(épinglés par $M$, $N$, $\sigma$) est donc une \uncertain{cat.} de Picard
stricte d'inv.\ $M$, $N$ — d'invariant $\in \mathrm{Ext}^2(S; M, N)$.

Il reste encore à voir à quelles conditions un triple $(M,N,\sigma)$
provient d'un $P$. \struck{(L'obstruction serait-elle dans
$\mathrm{Ext}^3(S;M,N)$ ?)}

C'est OK si $\sigma$ est la restriction d'une application bilin.\
(antisymétrique) $M \times M \to N$, $\varphi(x,y) = -\varphi(y,x)$ —
c'est le cas si $M_2$ \note{suit une demi-ligne biffée, illisible} est un
\uncertain{faisceau} « libre », p.ex.\ si $M$ est \uncertain{« libre »} sur
$\mathbf{Z}$. Cela \uncertain{montre} déjà que si $M$ est quotient d'un $L$
tel que le composé $L \xrightarrow{\ p\ } M \xrightarrow{\ \sigma\ } {}_2N$
est réalisable par une cat.\ de Picard (\uncertain{disons} $P$,
\uncertain{munie d'un élément de} $\mathrm{Ext}^2(S;L,N)$), la
« restriction » de cette $P$ : $(R,N)$ ($R = \mathrm{Ker}\,p$) définit une
classe dans $\mathrm{Ext}^2(S;R,N)$ déterminée mod.\ images d'éléments des
$\mathrm{Ext}^2(S;L,N)$, donc son image $\omega(\sigma)$ dans
$\mathrm{Ext}^3(S;M,N)$ est can.\ déterminée, et elle est nulle ssi
\uncertain{l'on peut}

\marginal{cette condition \uncertain{s'interprète} : $P$ (d'invariants
$(M,N,\sigma)$) admet une \struck{\ill} \uncertain{contrainte de comm.\ qui
en fasse un champ} de Picard \emph{strict}}

\page{16}
choisir $P$ (Pic.\ d'invariants $L$, $N$) telle que $P|_R$ soit splittée —
ou, ce qui revient au =, telle que $P$ provienne d'une cat.\ de Pic.\
d'inv.\ $M$, $N$. \note{le $N$ des invariants corrige un $M$} \struck{\ill}
\uncertain{On vérifie immédiatement} que $\omega(\sigma)$ ne dépend pas des
choix particuliers de $L \to M \to 0$, et que sa formation est
fonctorielle en $M$, $N$ \uncertain{et} compatible avec \struck{\ill}
\uncertain{les images} inverses par morphismes de topos. On se place dans
la situation universelle de $M \to {}_2N$, \uncertain{soit} $M
\xrightarrow{\ \sigma\ } {}_2N$ \uncertain{sur $B$} — \uncertain{aurions}
$\mathrm{Ext}^i(B; M,N) = 0$ pour $i > 0$, en particulier $\omega(\sigma)
\in \mathrm{Ext}^3(B; M,N) = 0$ est nulle, donc $\omega(\sigma)$ est nulle
\uncertain{dans tous les cas} !

\page{17}
\note{page 6 de la main de l'auteur}
$M$, $N$ faisceaux abéliens sur \uncertain{le} topos $X$. $\mathcal{P}$
2-catégorie des chps de Picard sur $X$ épinglés par \uncertain{$P$}.
\note{comprendre sans doute « épinglés par $(M,N)$ » ; la page ne porte
qu'une lettre} $\mathcal{P}$ est \uncertain{munie} de la composition de
Baer.

$\pi_0(\mathcal{P})$ = groupe des objets de $\mathcal{P}$ à équivalence
près :
\[
0 \longrightarrow \mathrm{Ext}^2(M,N) \longrightarrow \pi_0(\mathcal{P})
\longrightarrow \mathrm{Hom}(M, {}_2N) \longrightarrow 0
\]
\note{le premier argument du $\mathrm{Ext}^2$ se lit aussi bien $N$ ; la
suite exacte impose $M$}

$\pi_1(\mathcal{P})$ = groupe des classes d'\uncertain{isomorphie}
d'auto-équivalences d'un objet de $\mathcal{P}$ $= \mathrm{Ext}^1(M,N)$.

$\pi_0(\mathcal{P})$ = groupe des \uncertain{automorphismes du foncteur
identique} d'un objet $P$ de $\mathcal{P}$ $= \mathrm{Hom}(M,N)$.
\note{la page répète $\pi_0$ ; il s'agit du groupe suivant, le $\pi_2$}

$\mathcal{P}$ 2-catégorie de Picard \emph{stricte} ?

$P \otimes P$ = catégorie dont les objets sont $(L, L')$ avec
\uncertain{$\partial L \simeq L'$} \ill{}
\[
\mathrm{Hom}\bigl((L,L'), (U,U')\bigr)
= \bigl[\mathrm{Hom}(L,U) \times \mathrm{Hom}(L',U')\bigr] / N
\quad \text{(diagonale)}
\]

\section*{Splittages et la 2-gerbe associée à une Gr-catégorie (pages 18
et 19)}

\page{18}
\note{page 7 de la main de l'auteur. La page ouvre sur un tableau de
correspondances : deux blocs face à face reliés par une double flèche, un
bloc médian, et une flèche montant d'un quatrième bloc vers celui de
droite ; les lignes du haut du bloc de droite et du bloc médian sont en
partie biffées}

« La 3-catégorie des Gr-catégories $C$ avec $\pi_0(C) = G$, $\pi_1(C) =
M$ \note{une lettre biffée précède le $M$}, et 2-catégories des splittages
des Gr-catégories » ---
« "noyaux" et représentations des noyaux par extensions » --- « La
3-catégorie des 2-gerbes sur $B_G$ \struck{de faisceaux} \struck{de liens
\ill{}} de liens commutatifs \uncertain{donnés} (avec $G$ opérant
dessus) » --- « Liens sur $B_G$ et 2-gerbes des représentations des
liens ».

Soit $C$ une Gr-catégorie d'invariants $G$, $M$. On va lui associer une
2-gerbe \uncertain{$\widetilde{C}$} sur $B_G$, donc : à \uncertain{tout}
objet $X$ de $B_G$ (aux $G$-ens) on associe une 2-catégorie. Si $X$ est
l'objet final de $B_G$, ce sera la 2-catégorie des splittages de $C$.

\struck{Fixons pour \uncertain{tout} $x \in G$ un représentant $L(x) \in
\mathrm{Ob}\,C$, et pour tout couple $(x,y)$, d'un isom $\varphi_{x,y} :
L(x) \otimes L(y) \simeq L(xy)$ ; \ill{} $(L, \varphi)$ ($L : G \to
\mathrm{Ob}\,C$, $\varphi : G \times G \to \mathrm{Fl}\,C$) \ill{} d'une
2-cochaîne \ill{} $\lambda : G \times G \to M$ \ill{} $\psi_{x,y} =
\varphi_{x,y}\,\lambda(x,y)$, la donnée d'associativité \ill{} soit nulle.
L'ens.\ des $L$ splitt.\ \uncertain{forme un espace principal homogène}
sous $Z^2(G,M)$. Soient $(L,\varphi)$, $(L,\varphi')$ deux splittages
\uncertain{épinglés de type $G$, $M$}}
\note{tout ce passage est barré de grands traits obliques ; en marge et en
interligne, un renvoi : « En termes de $L$, $\varphi$ — $L : G \to
\mathrm{Ob}\,C$, $\varphi : G \times G \to \mathrm{Fl}(C)$ »}

Soit $C_0$ \uncertain{la} (Gr-)catégorie \uncertain{ci-décrite}, splittée
\struck{(i.e.\ \ill{} l'identité)} :
\[
\mathrm{Ob}\,C = G, \qquad
\mathrm{Fl}\,C = G \times M \rightrightarrows G, \qquad
s = b = \mathrm{pr}_1,
\]
\[
(g,m) \cdot (g,m') = (g, m+m'), \qquad
g \otimes g' = gg',
\]
\note{devant $g \otimes g' = gg'$, un premier départ
« $(g,m) \otimes (g',m')$ » est biffé}
\[
(g,m) \otimes (g',m') = (gg', m+m'), \qquad
a_{g,g',g''} = \mathrm{id}_{gg'g''}.
\]
\note{lire $C_0$ — la page écrit $C$}
\uncertain{Or} les splittages de $C$ \uncertain{sont les} équivalences des
Gr-catégories (épinglées) $S : C_0 \to C$, (\ill{}
\note{le $S$ porte un exposant — $(L,\varphi)$ ? — surchargé}

\page{19}
\struck{\ill} \uncertain{En tout cas}, les splittages de $C$ forment une
catégorie, qui est un 1-torseur sous la Gr-catégorie des auto-équivalences
de la Gr-catégorie $C_0$ avec elle-=. Quels sont les invariants de cette
dernière Gr-catégorie ?

Objets : les applications $\varphi : G \times G \to M$
(2-\uncertain{cochaînes}) qui sont des cobords.

Morphismes : $\varphi \to \psi$ est \uncertain{l'ens.} des 1-cochaînes
\uncertain{dont} le cobord est $\psi - \varphi$.

Objets à isom.\ près : $H^2(G,M)$ ; endomorphismes d'un objet :
$Z^1(G,M)$.
\note{la page porte bien « qui sont des cobords » quelques lignes avant
« à isom.\ près : $H^2(G,M)$ »}

\begin{tikzcd}[column sep=normal]
C_0 \arrow[r, bend left=30, "S"] \arrow[r, bend right=30, "S'"'] & C
\end{tikzcd}

\note{entre les deux flèches, une 2-flèche $\alpha$ ; la flèche $S$ porte
les données $L_x$, $\varphi_{x,y}$, la flèche $S'$ les données $L'_{x'}$,
$\varphi'_{x',y'}$}
\[
\alpha'(x) = \alpha(x)\,\varphi(x),
\]
\uncertain{d'où} $\delta\varphi = 0$ ;
\[
\mathrm{Hom}(\alpha, \alpha') = \{\, m \in M \mid \alpha'(x) =
\alpha(x)(gm - m) \,\}
\]
\note{entre les deux membres, plusieurs essais biffés — on y lit
« $\alpha' = \alpha(\delta g)$ » ; la fin du second est rognée par la
marge}

\note{un trait horizontal sépare ce qui précède de ce qui suit}

Soit $C$ une catégorie \struck{\ill} \uncertain{d'invariants $G$, $M$} ;
$X$ un ens.\ où $G$ opère. On va construire une 2-catégorie $S(C,X)$ :

0-objets : systèmes $L$, $\varphi$ \uncertain{ainsi} :

\begin{enumerate}
\item à \uncertain{tout} couple $(x,g)$ ($x \in X$, $g \in G$) un
\uncertain{$L(x,g) \in \mathrm{Ob}\,C$} ;
\item à \uncertain{tout} triple $(x,g,g')$ ($x \in X$, $g \in G$,
$g' \in G$ \note{l'appartenance du $g'$ est surchargée}) un isom
\[
L(x,g) \otimes L(gx, g') \simeq L(x, g'g).
\]
\note{le premier argument du second facteur est surchargé — un $gx$ récrit
sur autre chose}
\end{enumerate}

\marginal{$X \times G \rightrightarrows X$ ($\mathrm{pr}_1$ et
\uncertain{l'opération}) ; $x \xrightarrow{\ g\ } gx \xrightarrow{\ g'\ }
g'gx$}

\section*{Extensions d'une catégorie par un groupe (page 20)}

\page{20}
\note{la page ouvre une nouvelle numérotation — page 1 de la main de
l'auteur ; la note se poursuit au lot 2}
$\Gamma$ catégorie, $M$ groupe \struck{\ill}.

Une extension de $\Gamma$ par $M$ est la donnée :

a) pour toute flèche $u \in \mathrm{Fl}\,\Gamma$, d'un $M$-bitorseur $P_u$ ;

b) pour deux flèches composables $(u,v)$ de $\Gamma$, d'un isom
$\varphi_{u,v} : P_u \otimes P_v \simeq P_{uv}$, avec axiome de
commutativité du diagramme suivant, si $(u,v,w)$ composables (axiome
\uncertain{pentagonal}) :

\begin{tikzcd}[column sep=small, row sep=small, nodes={font=\scriptsize}]
(P_u \otimes P_v) \otimes P_w \arrow[r, "\varphi_{u,v} \otimes \mathrm{id}"]
\arrow[d, "\mathrm{can}"'] & P_{uv} \otimes P_w \arrow[dd, "\varphi_{uv,w}"] \\
P_u \otimes (P_v \otimes P_w) \arrow[d, "\mathrm{id} \otimes \varphi_{v,w}"'] & \\
P_u \otimes P_{vw} \arrow[r, "\varphi_{u,vw}"'] & P_{uvw}
\end{tikzcd}

\marginal{$\mathrm{Fl}_3(\Gamma)$, $\mathrm{Fl}_2(\Gamma)$,
$\mathrm{Fl}_1\Gamma$, $\mathrm{Fl}_0\Gamma$, chaque étage envoyé sur le
précédent par ses flèches de faces — quatre, trois, deux}

[NB — au lieu de se donner $M$, \uncertain{il suffit de se donner} une
catégorie \struck{de P\ill} \uncertain{$\widetilde{P}$} — \uncertain{ici la
catégorie des} $M$-bitorseurs.]

[Invariants : 1) $\mathrm{Fl}(\Gamma) \xrightarrow{\ \alpha\ }
\pi_0(\widetilde{P})$ ($=$ \uncertain{$\mathrm{Autext}(M)$}) ; 2)
$c \in H^2(\Gamma ; \widetilde{P})$ — (c'est un $H^2$ non commutatif, en
général sans élément neutre \struck{\ill} [commutatif]).]
\note{les deux numéros sont cerclés sur la page ; l'intérieur du $H^2$ est
griffonné puis récrit, un $\alpha$ y est biffé, et le $\widetilde{P}$ des
deux formules reste une lecture incertaine}
\marginal{$= \mathrm{Aut}(M)$ si $M$ \uncertain{commutatif}}

Une 2-extension de $\Gamma$ par $M$ est la donnée :

a) pour toute flèche $u \in \mathrm{Fl}\,\Gamma$, d'une $M$-gerbe $P_u$ ;

b) pour deux flèches composables $u$, $v$, d'une équivalence
$\varphi_{u,v} : P_u \otimes P_v \simeq P_{uv}$ ;

c) pour trois flèches composables, d'un isom \uncertain{rendant
commutatif} le diagramme pentagonal

\end{document}
